\documentclass[12pt,letterpaper]{article}

\usepackage{fontspec}
\usepackage{polyglossia}
\setdefaultlanguage{spanish}
\setmainfont{Liberation Serif}
\setmonofont{Liberation Mono}

\usepackage{geometry}
\geometry{margin=2.5cm}

\usepackage{xcolor}
\usepackage{listings}
\usepackage{booktabs}
\usepackage{enumitem}
\usepackage{ragged2e}
\usepackage{hyperref}

\definecolor{codegray}{gray}{0.95}
\definecolor{codegreen}{rgb}{0.0,0.45,0.0}
\definecolor{codeblue}{rgb}{0.1,0.1,0.6}

\lstset{
    backgroundcolor=\color{codegray},
    basicstyle=\ttfamily\small,
    keywordstyle=\color{codeblue}\bfseries,
    commentstyle=\color{codegreen}\itshape,
    stringstyle=\color{red!70!black},
    numbers=left,
    numberstyle=\tiny\color{gray},
    numbersep=8pt,
    breaklines=true,
    showstringspaces=false,
    frame=single,
    rulecolor=\color{gray!50},
    tabsize=4,
    extendedchars=true
}

\title{Pr\'actica 2. Productor -- Consumidor}
\author{Islas G\'omez Luis Antonio \and Ram\'irez Candia Alfredo \and P\'erez Flores Julio C\'esar \and Navarro Soto Mario Alberto}
\date{\today}

\begin{document}

%% ------------------------------------------------------------------
%% PORTADA
%% ------------------------------------------------------------------
\begin{titlepage}
    \centering
    \vspace*{0.5cm}

    {\Large\bfseries BENEM\'ERITA UNIVERSIDAD AUT\'ONOMA DE PUEBLA\par}
    \vspace{0.4cm}
    {\large Facultad de Ciencias de la Computaci\'on\par}
    \vspace{1.2cm}

    {\large Programaci\'on Concurrente y Paralela\par}

    \vfill

    {\Huge\bfseries Pr\'actica 2\par}
    \vspace{0.4cm}
    {\LARGE Productor -- Consumidor\par}
    \vspace{0.3cm}
    {\large (Primer parcial)\par}

    \vfill

    {\large\bfseries Equipo:\par}
    \vspace{0.3cm}
    Islas G\'omez Luis Antonio --- 202372860\par
    Ram\'irez Candia Alfredo --- 202375947\par
    P\'erez Flores Julio C\'esar --- 202247445\par
    Navarro Soto Mario Alberto --- 202264602\par

    \vfill

    {\large Docente: Hilda Castillo Zacatelco\par}
    \vspace{0.6cm}
    {\large \today\par}
\end{titlepage}

\tableofcontents
\newpage

%% ------------------------------------------------------------------
\section{Introducci\'on y objetivo}
El problema productor--consumidor consiste en dos procesos que comparten un
buffer: el productor deposita datos y el consumidor los retira. El objetivo de
esta pr\'actica es implementarlo \'unicamente con \texttt{fork()} y archivos de
texto, sin emplear tuber\'ias, se\~nales, \texttt{wait}, \texttt{exec},
sem\'aforos, mutex, monitores, memoria compartida ni colas de mensajes. La
sincronizaci\'on se resuelve con espera ocupada (ciclos vac\'ios) y retardos.

%% ------------------------------------------------------------------
\section{Marco te\'orico}
\begin{itemize}[leftmargin=1.4cm]
    \item \textbf{Regi\'on cr\'itica:} fragmento de c\'odigo que accede a un
    recurso compartido y no debe ejecutarse simult\'aneamente por dos procesos.
    \item \textbf{Condici\'on de carrera:} fallo producido cuando el resultado
    depende del orden no determinista en que los procesos acceden al recurso.
    \item \textbf{Exclusi\'on mutua:} garant\'ia de que solo un proceso est\'a en
    la regi\'on cr\'itica a la vez.
    \item \textbf{Espera ocupada:} el proceso que no puede avanzar repite un
    ciclo hasta que la condici\'on cambia; se a\~naden peque\~nos retardos
    (\texttt{usleep}) para ceder tiempo de CPU.
    \item \textbf{COBEGIN / COEND:} notaci\'on que encierra bloques de c\'odigo
    que pueden ejecutarse en paralelo.
\end{itemize}

%% ------------------------------------------------------------------
\section{Desarrollo}
\subsection{Estructura del programa}
Se crea un proceso productor con \texttt{fork()}. El proceso padre act\'ua como
consumidor y el hijo como productor. El buffer compartido es el archivo
\texttt{buffer.txt} (una l\'inea por elemento) y el candado es
\texttt{candado.txt}. El productor agrega al final del archivo y el consumidor
retira el primer elemento, por lo que se respeta el orden FIFO.

\subsection{Candado at\'omico}
La exclusi\'on mutua se logra creando el archivo candado con la bandera
\texttt{O\_CREAT | O\_EXCL}. El sistema operativo realiza esa creaci\'on de
forma indivisible: si el archivo ya existe, la llamada falla con
\texttt{EEXIST} y el proceso reintenta en un ciclo (espera ocupada). Al
terminar, el candado se libera con \texttt{unlink()}.

\begin{lstlisting}[language=C]
void adquirirCandado(){
    while(open(ARCHIVO_CANDADO, O_CREAT | O_EXCL | O_WRONLY, 0600) == -1){
        usleep(1000);
    }
}

void liberarCandado(){
    unlink(ARCHIVO_CANDADO);
}
\end{lstlisting}

\subsection{Regi\'on cr\'itica}
El conteo del buffer y su modificaci\'on se realizan \emph{dentro} del candado,
de modo que la verificaci\'on y la operaci\'on forman una sola acci\'on
at\'omica frente al otro proceso.

\begin{lstlisting}[language=C]
/* Productor */
adquirirCandado();
if(contarItems() < BUFFER_SIZE){
    agregarItem(numero);
    liberarCandado();
    /* ... */
} else {
    liberarCandado();
}
\end{lstlisting}

%% ------------------------------------------------------------------
\section{Problemas de sincronizaci\'on y condiciones de carrera}
Al compartir recursos mediante archivos surgen los siguientes riesgos, que la
soluci\'on presentada debe evitar:
\begin{enumerate}[leftmargin=1.4cm]
    \item \textbf{Condici\'on de carrera en el acceso concurrente.} Si los dos
    procesos leen y modifican el archivo sin coordinarse, sus operaciones se
    intercalan y el resultado depende del orden de ejecuci\'on.

    \item \textbf{Falta de atomicidad del candado.} Un candado construido
    leyendo y luego escribiendo el archivo deja una ventana entre ambas
    operaciones; ambos procesos podr\'ian encontrar el candado libre y entrar a
    la vez. Por eso se requiere una primitiva indivisible.

    \item \textbf{Verificaci\'on y modificaci\'on separadas.} Comprobar el
    conteo del buffer antes de acceder a \'el produce una condici\'on
    \emph{check-then-act}: entre la comprobaci\'on y la operaci\'on el buffer
    puede cambiar por el otro proceso.

    \item \textbf{P\'erdida de actualizaciones.} Como \texttt{quitarItem()}
    reescribe todo el archivo, un agregado que ocurra entre la lectura y la
    reescritura se perder\'ia. Igualmente, el productor podr\'ia desbordar el
    buffer o el consumidor extraer de un buffer vac\'io.
\end{enumerate}

%% ------------------------------------------------------------------
\section{Soluci\'on planteada}
La soluci\'on se basa en dos ideas: el candado se toma mediante una
operaci\'on at\'omica del sistema de archivos y toda verificaci\'on y
modificaci\'on del buffer se realiza dentro de la regi\'on cr\'itica. El
siguiente fragmento muestra el algoritmo implementado, encerrado entre
COBEGIN y COEND para representar los dos procesos que corren en paralelo.

\begin{lstlisting}[language=C]
// Variables compartidas: buffer.txt (buffer) y candado.txt (candado). 

// COBEGIN 
//Productor (proceso hijo) */
int numero = 1;
while(1){
    adquirirCandado();
    if(contarItems() < BUFFER_SIZE){
        agregarItem(numero);
        liberarCandado();
        printf("[Productor] genero: %d\n", numero);
        numero = (numero % 10) + 1;
    } else {
        liberarCandado();
    }
    usleep(200000);
}

// Consumidor (proceso padre) 
while(1){
    adquirirCandado();
    int item = -1;
    if(contarItems() > 0){
        item = quitarItem();
    }
    liberarCandado();
    if(item != -1){
        printf("\t\t[Consumidor] tomo: %d\n", item);
    }
    usleep(500000);
}
// COEND 
\end{lstlisting}

\noindent
Con esto, \texttt{adquirir}/\texttt{liberar} garantizan exclusi\'on mutua y las
verificaciones de lleno/vac\'io ya no compiten con el otro proceso.

%% ------------------------------------------------------------------
\section{Pruebas y resultados}
Se compil\'o con \texttt{gcc ComePica.c -o bin/comepica -Wall -Wextra} sin
errores y se ejecut\'o durante algunos segundos. La salida muestra que el
productor nunca supera los 10 elementos y el consumidor solo retira cuando hay
datos; el orden de retiro es FIFO.

\begin{lstlisting}
[Productor] genero: 1
[Productor] genero: 2
[Productor] genero: 3
        [Consumidor] tomo: 1
[Productor] genero: 4
...
        [Consumidor] tomo: 5
\end{lstlisting}

%% ------------------------------------------------------------------
\section{Conclusiones}
El uso de archivos de texto como recurso compartido es suficiente si se
construye un candado at\'omico y se define con precisi\'on la regi\'on
cr\'itica. La principal ense\~nanza es que no basta con ``bloquear'' el acceso:
separar la verificaci\'on de la operaci\'on deja abierta una condici\'on de
carrera. Colocar ambas acciones bajo el candado y usar una primitiva
indivisible del sistema de archivos elimina las p\'erdidas de actualizaciones y
respeta los l\'imites del buffer.

\end{document}
